\chapter{Ukuran Hasil Kali, Fubini, Penggantian Variabel}
\label{ch:b3:product}

Teori Lebesgue satu dimensi menjadi kalkulus berdimensi banyak
lewat dua teorema. \emph{Tonelli--Fubini} mengatakan bahwa
integral atas hasil kali adalah integral berulang --- sehingga pengirisannya
sah, dalam urutan mana pun, di bawah hipotesis yang sungguh dapat
diperiksa. Sedangkan \emph{rumus penggantian variabel} mengangkut
integral sepanjang difeomorfisma $\mathcal C^1$, dengan
determinan Jacobi sebagai kurs tukar volumenya; dan kita membuktikannya
secara \omterm{def:b3:complete:complete}{lengkap}, mulai dari kasus linearnya, tempat ia menerangkan
determinan itu \emph{apa}. Terapannya lalu beruntun: rumus kue
berlapis, konvolusi, \omterm{ex:b3:product:polar}{koordinat kutub}, volume
bola-$n$ --- dan, pada soal akhir pekannya, rumus Stirling
beserta analisis galat yang jujur.

\section{Aljabar-$\sigma$ hasil kali dan ukuran hasil kali}

\begin{definition}\label{def:b3:product:sigma}
Untuk ruang \omterm{def:b3:lebesgue:measurable}{terukur} $(X, \mathcal A)$ dan $(Y, \mathcal B)$,
\emph{aljabar-$\sigma$ hasil kali}\index{aljabar-$\sigma$
hasil kali} $\mathcal A \otimes \mathcal B$ pada $X \times
Y$ dibangkitkan oleh \emph{persegi panjang} $A \times B$ (dengan $A \in
\mathcal A$ dan $B \in \mathcal B$) --- yakni sebuah sistem-$\pi$. Untuk $E
\subseteq X\times Y$ dan $x \in X$, \emph{penampangnya} adalah $E_x
= \{y : (x,y) \in E\}$; sedangkan untuk sebuah fungsi $f$ pada hasil kalinya,
$f_x = f(x, \cdot)$.
\end{definition}

\begin{proposition}\label{prop:b3:product:sections}
(a) Jika $E \in \mathcal A\otimes\mathcal B$, maka setiap penampang $E_x
\in \mathcal B$ (dan setangkup dengan itu); dan jika $f$ \omterm{def:b3:lebesgue:measurable}{terukur} terhadap $\mathcal
A\otimes\mathcal B$, maka setiap $f_x$ \omterm{def:b3:lebesgue:measurable}{terukur} terhadap $\mathcal
B$.
(b) $\mathcal B(\R^m)\otimes\mathcal B(\R^n) = \mathcal
B(\R^{m+n})$.
\end{proposition}

\begin{proof}
(a) Himpunan baiknya: $\{E : E_x \in \mathcal B\ \forall x\}$ merupakan
aljabar-$\sigma$ (sebab penampangnya komutatif dengan komplemen dan
gabungan terbilang) yang memuat persegi panjangnya. Untuk $f$:
$(f_x)^{-1}(B) = (f^{-1}(B))_x$.
(b) ($\subseteq$) Untuk persegi panjang himpunan Borel: cukup bahwa
kotak terbuka$\times$terbuka bersifat Borel di $\R^{m+n}$ (sebab kotaknya terbuka)
dan bahwa persegi panjang Borel umumnya berupa limit --- lewat himpunan baik
lagi: $\{A : A\times\R^n \in \mathcal B(\R^{m+n})\}$ merupakan
aljabar-$\sigma$ yang memuat \omterm{def:b3:topology:topology}{himpunan terbukanya}; lalu iriskan keduanya.
($\supseteq$) Setiap \omterm{def:b3:topology:topology}{himpunan terbuka} $\R^{m+n}$ berupa gabungan terbilang
kotak terbuka rasional $U\times V$: yang termuat di aljabar-$\sigma$
hasil kalinya.
\end{proof}

\begin{theorem}[Ukuran hasil kali]\label{thm:b3:product:existence}
Misalkan $(X, \mathcal A, \mu)$ dan $(Y, \mathcal B, \nu)$ bersifat
\emph{berhingga-$\sigma$}. Untuk setiap $E \in \mathcal
A\otimes\mathcal B$, fungsi $x \mapsto \nu(E_x)$ bersifat
\omterm{def:b3:lebesgue:measurable}{terukur}, dan\index{ukuran hasil kali}
\[
(\mu\otimes\nu)(E) = \int_X \nu(E_x)\,\dd\mu(x)
\]
mendefinisikan satu-satunya \omterm{def:b3:measure:measure}{ukuran} pada $\mathcal A\otimes\mathcal B$
dengan $(\mu\otimes\nu)(A\times B) = \mu(A)\nu(B)$. \omterm{def:b3:measure:measure}{Ukuran} itu
berhingga-$\sigma$ dan setangkup: sebab \omterm{def:b3:measure:measure}{ukuran} yang sama diperoleh dengan
mengintegralkan penampang-$x$ terhadap $\nu$.
\end{theorem}

\begin{proof}
\emph{\omterm{def:b3:lebesgue:measurable}{Keterukuran} $x \mapsto \nu(E_x)$.} Mula-mula misalkan $\nu$
berhingga. Kelas $\mathcal D$ berisi $E$ yang pemetaannya \omterm{def:b3:lebesgue:measurable}{terukur}
memuat persegi panjangnya (sebab $\nu((A\times B)_x) =
\nu(B)\mathbf 1_A(x)$) dan merupakan sistem-$\lambda$: sebab untuk $E
\subseteq F$ di $\mathcal D$, $\nu((F\setminus E)_x) =
\nu(F_x) - \nu(E_x)$ (menurut keberhinggaannya); sedangkan untuk $E_n \uparrow E$,
$\nu((E_n)_x) \uparrow \nu(E_x)$ (menurut \omterm{def:b3:topology:continuity}{kekontinuan} dari bawah), dan
limit monoton \omterm{def:b3:lebesgue:measurable}{fungsi terukur} tetap \omterm{def:b3:lebesgue:measurable}{terukur}.
Persegi panjangnya membentuk sistem-$\pi$: sehingga Dynkin
(\cref{thm:b3:measure:dynkin}) memberikan $\mathcal D = \mathcal
A\otimes\mathcal B$. Jika $\nu$ berhingga-$\sigma$, tulislah $Y =
\bigcup Y_k$ dengan $Y_k \uparrow$ dan $\nu(Y_k) < \infty$: maka $\nu(E_x) =
\lim_k\nu_k(E_x)$ dengan $\nu_k = \nu(\cdot\cap Y_k)$ yang berhingga.

\emph{\omterm{def:b3:measure:measure}{Ukurannya}.} Keaditifan-$\sigma$ dari $E \mapsto
\int\nu(E_x)\dd\mu$ menyusul dari
\cref{cor:b3:lebesgue:additivity} (sebab penampang himpunan yang saling lepas
tetap saling lepas). Dan pada persegi panjangnya ia memberikan $\mu(A)\nu(B)$.
\emph{Ketunggalannya}: dua calonnya bersesuaian pada sistem-$\pi$ berisi
persegi panjang; sedangkan keberhinggaan-$\sigma$-nya menyediakan persegi panjang $X_k\times
Y_k \uparrow X\times Y$ yang berukuran berhingga:
\cref{thm:b3:measure:uniqueness}. Kesetangkupannya: sebab
konstruksi urutan lainnya juga merupakan \omterm{def:b3:measure:measure}{ukuran} yang bersesuaian pada
persegi panjangnya --- yang tunggal, jadi sama.
\end{proof}

\begin{definition}\label{def:b3:product:lebesgue}
\emph{Ukuran Lebesgue pada $\R^d$}\index{ukuran Lebesgue pada
$\R^d$} adalah $\lambda_d = \lambda\otimes\cdots\otimes\lambda$
(dengan $d$ faktor; sedangkan keasosiatifan konstruksinya diperiksa pada
kotak lalu dirambatkan lewat ketunggalannya). Ia satu-satunya \omterm{def:b3:measure:measure}{ukuran}
Borel yang memberi setiap kotak $\prod\intoc{a_i}{b_i}$ volumenya
$\prod(b_i - a_i)$; ia invarian terhadap translasi (sebab translasinya
bersesuaian pada kotak), berhingga-$\sigma$, dan \omterm{def:b3:complete:complete}{lengkap} setelah
\omterm{thm:b3:complete:completion}{pelengkapan} Carathéodory --- lalu kita menulis $\lambda_d$ untuk
\omterm{def:b3:measure:measure}{ukuran} yang telah dilengkapkan itu dan mengintegralkan seturut itu.
\end{definition}

\section{Tonelli dan Fubini}

\begin{theorem}[Tonelli]\label{thm:b3:product:tonelli}
Misalkan $\mu, \nu$ berhingga-$\sigma$ dan $f \colon X\times Y \to [0,
+\infty]$ \omterm{def:b3:lebesgue:measurable}{terukur}. Maka $x \mapsto \int_Y f_x\,\dd\nu$ bersifat
\omterm{def:b3:lebesgue:measurable}{terukur} dan\index{teorema Tonelli}
\[
\int_{X\times Y}f\,\dd(\mu\otimes\nu)
= \int_X\Bigl(\int_Y f(x,y)\,\dd\nu(y)\Bigr)\dd\mu(x)
= \int_Y\Bigl(\int_X f(x,y)\,\dd\mu(x)\Bigr)\dd\nu(y).
\]
\end{theorem}

\begin{proof}
Lewat mesin bakunya. Untuk $f = \mathbf 1_E$ ini adalah
\cref{thm:b3:product:existence} (beserta bentuk setangkupnya). Lewat
kelinearannya ia berlaku bagi $f \geq 0$ yang sederhana. Sedangkan untuk $f \geq
0$ yang umum: ambillah $s_n \nearrow f$ yang sederhana
(\cref{thm:b3:lebesgue:approximation}); maka $\int_Y(s_n)_x
\dd\nu \nearrow \int_Y f_x\dd\nu$ untuk setiap $x$ (lewat kekonvergenan monoton di $Y$),
sehingga anggota kirinya konvergen lewat kekonvergenan monoton di $X$, sedangkan
$\int s_n\,\dd(\mu\otimes\nu) \nearrow \int f$ lewat kekonvergenan monoton pada
hasil kalinya.
\end{proof}

\begin{theorem}[Fubini]\label{thm:b3:product:fubini}
Misalkan $\mu, \nu$ berhingga-$\sigma$ dan $f \in L^1(\mu\otimes\nu)$. Maka
untuk $\mu$-hampir setiap $x$ penampang $f_x$ \omterm{def:b3:lebesgue:l1}{terintegralkan} terhadap $\nu$, lalu
fungsi $x \mapsto \int f_x\dd\nu$ yang terdefinisi hampir di mana-mana bersifat
\omterm{def:b3:lebesgue:l1}{terintegralkan}, dan kedua integral berulangnya sama-sama bernilai
$\int f\,\dd(\mu\otimes\nu)$.\index{teorema Fubini}
\end{theorem}

\begin{proof}
Tonelli yang diterapkan pada $\abs f$ menunjukkan bahwa $\varphi(x) = \int\abs{f_x}
\dd\nu$ berintegral berhingga, jadi berhingga hampir di mana-mana: sehingga $f_x \in
L^1(\nu)$ untuk hampir setiap $x$. Pecahlah $f = f^+ - f^-$ (pada kasus realnya;
sedangkan kasus kompleksnya menurut komponennya): maka Tonelli menghitung setiap integral
berulang $f^\pm$ sebagai $\int f^\pm\dd(\mu\otimes\nu) <
\infty$, dan selisih yang terdefinisi hampir di mana-mana itu berintegral
selisihnya. Setangkup dengan itu untuk urutan yang lain.
\end{proof}

\begin{method}\label{met:b3:product:fubini}
Untuk menukar dua integral (atau sebuah integral dengan sebuah jumlah, atau dua
jumlah): jika integrannya \emph{tak negatif}, tukarlah
sebebasnya (lewat Tonelli). Selain itu, terapkanlah dahulu Tonelli pada $\abs f$
dalam urutan mana pun yang lebih mudah ditaksir; dan jika hasilnya
berhingga, maka Fubini mengesahkan penukarannya. Jangan pernah melewati
pemeriksaan $\abs f$ itu: sebab integran \cref{exo:b3:product:4} mempunyai dua
integral berulang \emph{bernilai berbeda}.
\end{method}

\begin{proposition}[Kue berlapis]\label{prop:b3:product:layercake}
Untuk $f \geq 0$ yang \omterm{def:b3:lebesgue:measurable}{terukur} pada $(X, \mathcal A, \mu)$ yang
berhingga-$\sigma$:\index{rumus kue berlapis}
\[
\int_X f\,\dd\mu = \int_0^{+\infty}\mu(\{f > t\})\,\dd t,
\qquad
\int_X f^p\,\dd\mu = p\int_0^{+\infty}t^{p-1}\mu(\{f >
t\})\,\dd t \quad (p \geq 1).
\]
\end{proposition}

\begin{proof}
Terapkan Tonelli pada $\mathbf 1_{\{(x,t) : 0 < t < f(x)\}}$ di $X
\times \intoo0{+\infty}$ (yang \omterm{def:b3:lebesgue:measurable}{terukur}: sebab ia $\{(x,t): f(x) - t
> 0\}\cap\{t > 0\}$, yakni gabungan bertipe Borel atas
$(x,t)\mapsto f(x) - t$ yang \omterm{def:b3:lebesgue:measurable}{terukur}): maka mengintegralkan terhadap $t$ lebih dahulu
memberikan $\int f\,\dd\mu$; sedangkan terhadap $x$ lebih dahulu, $\int_0^\infty\mu(f >
t)\dd t$. Untuk $f^p$: substitusikan $t = s^p$ pada $\int\mu(f^p > t)
\dd t$, yakni terapkan rumus pertamanya pada $f^p$ lalu gantilah
variabelnya pada integral satu dimensinya (sebab $\{f^p > s^p\} =
\{f > s\}$).
\end{proof}

\begin{theorem}[Konvolusi pada $L^1$]\label{thm:b3:product:convolution}
Untuk $f, g \in L^1(\R^d, \lambda_d)$, integral
\[
(f * g)(x) = \int_{\R^d} f(x - y)\,g(y)\,\dd y
\]
konvergen mutlak untuk hampir setiap $x$, mendefinisikan $f * g \in
L^1(\R^d)$ dengan $\norm{f*g}_1 \leq \norm f_1\norm g_1$, dan
$*$ bersifat komutatif serta asosiatif.\index{konvolusi}
\end{theorem}

\begin{proof}
Pemetaan $(x, y) \mapsto f(x-y)g(y)$ \omterm{def:b3:lebesgue:measurable}{terukur} (sebab $(x,y)\mapsto x -
y$ \omterm{def:b3:topology:continuity}{kontinu}; lalu komposisikan dan kalikan). Lewat Tonelli:
\[
\int\!\!\int \abs{f(x-y)}\abs{g(y)}\,\dd y\,\dd x
= \int\abs{g(y)}\Bigl(\int\abs{f(x - y)}\dd x\Bigr)\dd y
= \norm f_1\norm g_1 < \infty
\]
(lewat keinvarianan translasi $\lambda_d$ pada integral dalamnya).
Jadi integral rangkapnya berhingga; lalu Fubini memberikan kekonvergenan
mutlak hampir di mana-mana beserta batas normanya $\norm{f*g}_1 \leq \norm
f_1\norm g_1$. Kekomutatifannya: substitusikan $y \mapsto x - y$
(lewat keinvarianan translasi dan pencerminan --- sebab keinvarianan
pencerminannya berlaku pada kotak, jadi di mana-mana lewat ketunggalannya).
Keasosiatifannya: lewat Tonelli--Fubini pada integral rangkap tiga.
\end{proof}

\section{Penggantian variabel}

\begin{theorem}[Penggantian variabel linear]
\label{thm:b3:product:linearchange}
Untuk $T \in GL_d(\R)$ dan $A \in \mathcal B(\R^d)$:
$\lambda_d(T(A)) = \abs{\det T}\,\lambda_d(A)$; sehingga
$\int f(y)\dd y = \abs{\det T}\int f(Tx)\,\dd x$ untuk $f \geq 0$
atau yang \omterm{def:b3:lebesgue:l1}{terintegralkan}.\index{penggantian variabel linear}
\end{theorem}

\begin{proof}
\omterm{def:b3:measure:measure}{Ukuran} $\mu_T(A) = \lambda_d(T(A))$ merupakan \omterm{def:b3:measure:measure}{ukuran} Borel
(sebab \omterm{def:b3:topology:continuity}{homeomorfisma} mengawetkan himpunan Borel,
\cref{pb:b3:measure:1}), invarian terhadap translasi
(sebab $T(A + x) = T(A) + Tx$), dan berhingga pada kotak satuannya: sehingga lewat
pencirian \omterm{def:b3:measure:lebesgueouter}{ukuran Lebesgue}
(\cref{exo:b3:measure:6}, yang buktinya berlaku kata demi kata di
$\R^d$ dengan kubus diadik), $\mu_T = c(T)\lambda_d$ dengan $c(T)
= \lambda_d(T(\intco01^d))$. Pemetaan $T \mapsto c(T)$ bersifat
multiplikatif (sebab $c(ST) = c(S)c(T)$, lewat pengomposisian), sehingga
cukup dihitung $c$ pada pembangkit $GL_d$: yakni
matriks elementernya. Untuk yang diagonal $\operatorname{diag}(a, 1, \dots, 1)$:
ia memetakan kubus satuannya ke kotak bervolume $\abs a$: jadi $c = \abs a =
\abs\det$. Untuk transposisi koordinatnya: ia mempermutasikan kubusnya: jadi $c
= 1 = \abs\det$. Untuk transveksi $T(x) = x + \alpha
x_2e_1$: peta kubus satuannya adalah prisma yang tergeser; lalu lewat
Tonelli \omterm{def:b3:measure:measure}{ukurannya} adalah $\int\lambda_1(\text{penampang})\dd x_2
\cdots \dd x_d = 1$, sebab setiap penampang-$x_1$-nya berupa selang
berpanjang $1$: jadi $c = 1 = \abs{\det}$. Setiap matriks terbalikkan merupakan
hasil kali matriks semacam itu (lewat eliminasi Gauss), dan baik $c$ maupun
$\abs\det$ bersifat multiplikatif: sehingga $c(T) = \abs{\det T}$. Lalu
rumus integralnya menyusul lewat mesin bakunya (indikator,
sederhana, kekonvergenan monoton).
\end{proof}

\begin{theorem}[Penggantian variabel]\label{thm:b3:product:changeofvar}
Misalkan $U, V \subseteq \R^d$ terbuka dan $\Phi \colon U \to V$ sebuah
difeomorfisma $\mathcal C^1$. Untuk setiap $f \colon V
\to [0, +\infty]$ yang \omterm{def:b3:lebesgue:measurable}{terukur} (atau $f \in L^1(V)$):\index{rumus penggantian
variabel}
\[
\int_V f(y)\,\dd y = \int_U
f\bigl(\Phi(x)\bigr)\,\abs{\det D\Phi(x)}\,\dd x .
\]
\end{theorem}

\begin{proof}
Tulis $J(x) = \abs{\det D\Phi(x)}$. Jantung buktinya adalah
ketaksamaan
\[
\lambda_d\bigl(\Phi(A)\bigr) \leq \int_A J\,\dd\lambda_d
\qquad\text{untuk setiap Borel } A \subseteq U;
\tag{$*$}
\]
lalu Langkah 4 di bawah meningkatkan $(*)$ --- yang diterapkan pada $\Phi$ maupun
$\Phi^{-1}$ --- menjadi kesamaan teoremanya. Perhatikan bahwa
$\Phi^{-1}$ sendiri merupakan difeomorfisma $\mathcal C^1$ dengan
Jacobi $\abs{\det D\Phi^{-1}(y)} = J(\Phi^{-1}y)^{-1}$
(lewat aturan rantai pada $\Phi\circ\Phi^{-1} = \mathrm{id}$).

\emph{Langkah 1: $(*)$ untuk kubus, dengan faktor kesesatannya.} Tetapkan sebuah
kubus tertutup $Q \subseteq U$ berpusat $x_0$ dan bersisi $2r$
(yakni bola norma-sup). Klaimnya: untuk setiap $\varepsilon > 0$, jika $\Phi$
dapat diturunkan pada $Q$ dengan
$\norm{D\Phi(x) - D\Phi(x_0)} \leq \varepsilon$ pada $Q$
(dalam \omterm{def:b3:banach:operator}{norma operator} untuk norma-sup), maka
\[
\Phi(Q) \subseteq \Phi(x_0) + D\Phi(x_0)\Bigl(\,\bigl(1 +
\varepsilon\norm{D\Phi(x_0)^{-1}}\bigr)\,(Q - x_0)\Bigr),
\]
sebab untuk $x \in Q$, ketaksamaan nilai rata-rata yang diterapkan pada
$\Phi(x) - \Phi(x_0) - D\Phi(x_0)(x - x_0)$ memberikan
$\norm{\Phi(x) - \Phi(x_0) - D\Phi(x_0)(x-x_0)}_\infty \leq
\varepsilon\norm{x - x_0}_\infty \leq \varepsilon r$, dan
$D\Phi(x_0)^{-1}$ menarik cacat itu menjadi
pembesaran kubusnya sebesar $\varepsilon\norm{D\Phi(x_0)^{-1}}\,r$.
Menurut \cref{thm:b3:product:linearchange},
\[
\lambda_d(\Phi(Q)) \leq \abs{\det D\Phi(x_0)}\,
\bigl(1 + \varepsilon\,C\bigr)^{d}\,\lambda_d(Q),
\qquad C = \sup_{Q}\norm{D\Phi(\cdot)^{-1}} .
\]

\emph{Langkah 2: $(*)$ untuk kubus \omterm{def:b3:topology:compact}{kompak}, lewat pembagian.} Misalkan $Q
\subseteq U$ sebuah kubus \omterm{def:b3:topology:compact}{kompak} dan $\varepsilon > 0$. Pada $Q$,
$D\Phi$ \omterm{def:b3:topology:continuity}{kontinu} seragam dan $\norm{D\Phi^{-1}}$
terbatas (menurut kekompakannya); lalu bagilah $Q$ menjadi $2^{kd}$ subkubus
$Q_i$ yang cukup kecil sehingga osilasi $D\Phi$ pada masing-masingnya
$\leq \varepsilon$. Langkah 1 pada setiap subkubusnya (berpusat $x_i$):
\[
\lambda_d(\Phi(Q)) \leq \sum_i\lambda_d(\Phi(Q_i))
\leq (1 + C\varepsilon)^d \sum_i \abs{\det
D\Phi(x_i)}\,\lambda_d(Q_i)
\leq (1 + C\varepsilon)^d\Bigl(\int_Q J + \varepsilon'\Bigr),
\]
dengan langkah terakhirnya sebab $\sum_i\abs{\det D\Phi(x_i)}\mathbf
1_{Q_i} \to J$ secara seragam pada $Q$ (menurut \omterm{def:b3:topology:continuity}{kekontinuan} $\det D\Phi$)
--- yakni pembandingan jumlah Riemann. Lalu biarkan $\varepsilon \to 0$: sehingga $(*)$
berlaku bagi kubus \omterm{def:b3:topology:compact}{kompak}.

\emph{Langkah 3: $(*)$ untuk setiap Borel $A$.} Fungsi himpunan $A
\mapsto \lambda_d(\Phi(A))$, pada himpunan bagian Borel dari $U$, merupakan
\omterm{def:b3:measure:measure}{ukuran} (sebab $\Phi$ bijeksi ke $V$ yang mengawetkan himpunan Borel
dan kelepasan terbilang), dan demikian pula $A \mapsto \int_AJ$.
Setiap himpunan bagian terbuka $U$ berupa gabungan terbilang
kubus \omterm{def:b3:topology:compact}{kompak} diadik yang hampir saling lepas (lewat penguraian diadik
bakunya: ambillah kubus diadik maksimal yang termuat di \omterm{def:b3:topology:topology}{himpunan
terbukanya}), dan kedua \omterm{def:b3:measure:measure}{ukurannya} aditif melintasinya (sebab batas
kubusnya bersifat nol-$\lambda_d$, dan peta-$\Phi$-nya bersifat
nol menurut Langkah 2 yang diterapkan pada selimut kubus tipis atas sisinya):
sehingga $(*)$ beralih dari kubus ke \omterm{def:b3:topology:topology}{himpunan terbuka}. Untuk Borel $A$ yang umum:
habiskan $U$ oleh \omterm{def:b3:topology:compact}{kompak} $K_m \uparrow U$ dengan $K_m \subseteq
\mathring K_{m+1}$, lalu tetapkan $m$; maka pada $\mathring K_{m+1}$, $J$
terbatas oleh suatu $M_m$. Lewat keteraturan luar $\lambda_d$
(dengan bukti seperti pada \cref{thm:b3:measure:regularity}, memakai kotak),
pilihlah \omterm{def:b3:topology:topology}{himpunan terbuka} $O_n$ dengan $A \cap K_m \subseteq O_n
\subseteq \mathring K_{m+1}$ dan $\lambda_d\bigl(O_n \setminus
(A\cap K_m)\bigr) \to 0$. Maka
\[
\lambda_d\bigl(\Phi(A\cap K_m)\bigr) \leq
\lambda_d\bigl(\Phi(O_n)\bigr) \leq \int_{O_n}J
\leq \int_{A\cap K_m}J + M_m\,\lambda_d\bigl(O_n\setminus(A\cap
K_m)\bigr) \xrightarrow[n\to\infty]{} \int_{A\cap K_m}J .
\]
Lalu biarkan $m \to \infty$: \omterm{def:b3:topology:continuity}{kekontinuan} dari bawah pada kirinya, dan kekonvergenan monoton
pada kanannya. Dengan demikian $(*)$ tegak.

\emph{Langkah 4: kesamaannya dan rumus integralnya.} Mula-mula
perluaslah $(*)$ dari himpunan ke integral: untuk setiap $g
\geq 0$ yang \omterm{def:b3:lebesgue:measurable}{terukur} pada $V$,
\[
\int_V g(y)\,\dd y \leq \int_U g(\Phi(x))\,J(x)\,\dd x .
\tag{$**$}
\]
Sungguh, untuk $g = \mathbf 1_B$ ini adalah $(*)$ dengan $A =
\Phi^{-1}(B)$; lalu kelinearannya memperluasnya ke $g$ yang sederhana, dan kekonvergenan monotonnya ke
setiap $g \geq 0$ (yakni mesin bakunya). Kini terapkan $(**)$ dua kali:
mula-mula pada $g$, lalu --- untuk difeomorfisma $\Phi^{-1}$ ---
pada fungsi $x \mapsto g(\Phi(x))J(x)$:
\[
\int_Vg \leq \int_U g(\Phi(x))J(x)\dd x
\leq \int_V g(y)\,J(\Phi^{-1}y)\,\abs{\det
D\Phi^{-1}(y)}\,\dd y = \int_V g ,
\]
sebab $J(\Phi^{-1}y)\abs{\det D\Phi^{-1}(y)} = \abs{\det\bigl(
D\Phi(\Phi^{-1}y)\,D\Phi^{-1}(y)\bigr)} = 1$ (lewat aturan rantai pada
$\Phi\circ\Phi^{-1} = \mathrm{id}$). Jadi semua ketaksamaannya menjadi
kesamaan: sehingga rumusnya berlaku bagi $g \geq 0$, dan bagi fungsi
$L^1$ lewat penguraiannya.
\end{proof}

\begin{example}[Koordinat kutub; fungsi Gauss lagi]
\label{ex:b3:product:polar}
Pemetaan $\Phi(r, \theta) = (r\cos\theta, r\sin\theta)$ merupakan difeomorfisma
$\mathcal C^1$ dari $\intoo0{+\infty}\times\intoo0{2\pi}$
ke $\R^2$ dikurangi sebuah setengah garis (yakni himpunan nol), dengan $\det D\Phi =
r$:\index{koordinat kutub}
\[
\int_{\R^2}f(x, y)\,\dd x\,\dd y
= \int_0^{2\pi}\!\!\int_0^{+\infty}
f(r\cos\theta, r\sin\theta)\,r\,\dd r\,\dd\theta .
\]
Untuk $f = \eu^{-x^2-y^2}$, Tonelli dan rumus ini memberikan
\[
G^2 = \Bigl(\int_\R \eu^{-x^2}\dd x\Bigr)^2
= \int_{\R^2}\eu^{-x^2-y^2}
= 2\pi\int_0^\infty r\eu^{-r^2}\dd r = \pi:
\]
yakni bukti klasik dua baris bagi $G = \sqrt\pi$, yang kini sepenuhnya
dibenarkan (bandingkan dengan bukti parameternya pada
\cref{pb:b3:lebesgue:1}).
\end{example}

\begin{theorem}[Volume bola satuan]
\label{thm:b3:product:ballvolume}
Misalkan $v_d = \lambda_d(B(0,1))$ di $\R^d$. Maka\index{volume
bola satuan}
\[
v_d = \frac{\pi^{d/2}}{\Gamma\bigl(\frac d2 + 1\bigr)} :
\qquad
v_1 = 2,\quad v_2 = \pi,\quad v_3 = \tfrac{4\pi}3,\quad
v_4 = \tfrac{\pi^2}2,\ \dots
\]
\end{theorem}

\begin{proof}
Hitunglah $I = \int_{\R^d}\eu^{-\norm x_2^2}\dd\lambda_d$ dua kali.
Lewat Tonelli ia terfaktorkan: $I = G^d = \pi^{d/2}$. Sedangkan lewat rumus kue
berlapis (\cref{prop:b3:product:layercake}) dengan $f = \eu^{-
\norm x^2}$, yang himpunan arasnya berupa bola: $\{f > t\} =
B\bigl(0, \sqrt{-\ln t}\bigr)$ untuk $0 < t < 1$, yang berukuran
$v_d(-\ln t)^{d/2}$ (sebab pendilatan oleh $\rho$ menskalakan $\lambda_d$ dengan
$\rho^d$: \cref{thm:b3:product:linearchange}), sehingga
\[
I = \int_0^1 v_d\,(-\ln t)^{d/2}\,\dd t
\overset{t = \eu^{-s}}{=}
v_d\int_0^\infty s^{d/2}\eu^{-s}\,\dd s
= v_d\,\Gamma\Bigl(\frac d2 + 1\Bigr).
\]
Lalu samakan keduanya. (Nilainya: $\Gamma(\frac32) = \frac{\sqrt\pi}2$,
$\Gamma(2) = 1$, dan seterusnya.) Perhatikan bahwa $v_d \to 0$ saat $d \to \infty$ ---
dan soal akhir pekannya mengukur secepat apa, lewat Stirling.
\end{proof}

\section{Latihan}

\begin{exercise}[$\star$]\label{exo:b3:product:1}
Misalkan $\mu$ \omterm{def:b3:measure:measure}{ukuran} cacah pada $(\intcc01, \mathcal
B(\intcc01))$ (yang tak berhingga-$\sigma$) dan $\lambda$ \omterm{def:b3:measure:lebesgueouter}{ukuran
Lebesgue}, lalu misalkan $\Delta = \{(x,x)\}$ diagonal di
$\intcc01^2$. Tunjukkan bahwa $\Delta$ \omterm{def:b3:lebesgue:measurable}{terukur}, lalu hitunglah
kedua integral berulang $\mathbf 1_\Delta$ terhadap
$\lambda$ dan $\mu$: keduanya berbeda. Hipotesis mana pada
\cref{thm:b3:product:tonelli} yang gagal?
\end{exercise}

\begin{exercise}[$\star$]\label{exo:b3:product:2}
Benarkan penukarannya lalu turunkan kembali integral Dirichlet: untuk
$A > 0$,
\[
\int_0^A\frac{\sin x}x\,\dd x
= \int_0^A\!\!\int_0^{+\infty}\eu^{-xy}\sin x\,\dd y\,\dd x
= \int_0^{+\infty}\!\!\int_0^A \eu^{-xy}\sin x\,\dd x\,\dd y,
\]
hitunglah integral dalamnya dalam bentuk tertutup, lalu biarkan $A \to
+\infty$ (dengan mendominasi integral-$y$-nya) untuk memperoleh
$\int_0^\infty\frac{\sin x}x\dd x = \frac\pi2$.
\end{exercise}

\begin{exercise}[$\star\star$]\label{exo:b3:product:3}
(a) Buktikan bahwa untuk $f \geq 0$ yang \omterm{def:b3:lebesgue:measurable}{terukur} dan $\mu$ yang berhingga:
$\sum_{n\geq1}\mu(\{f \geq n\}) \leq \int f\,\dd\mu \leq
\mu(X) + \sum_{n\geq1}\mu(\{f\geq n\})$: sehingga keterintegralannya adalah
keterjumlahan \omterm{def:b3:measure:measure}{ukuran} ekornya.
(b) Turunkan bahwa $f \in L^1(\mu)$ (dengan $\mu$ berhingga) jika dan hanya jika
$\sum_n\mu(\abs f \geq n) < \infty$.
\end{exercise}

\begin{exercise}[$\star\star$]\label{exo:b3:product:4}
Untuk $f(x, y) = \dfrac{x^2 - y^2}{(x^2 + y^2)^2}$ pada
$\intoo01^2$, tunjukkan bahwa
\[
\int_0^1\!\!\int_0^1 f\,\dd y\,\dd x = \frac\pi4,
\qquad
\int_0^1\!\!\int_0^1 f\,\dd x\,\dd y = -\frac\pi4
\]
\emph{(perhatikan $f = \partial_y\bigl(\frac{y}{x^2+y^2}\bigr)$)},
lalu periksalah langsung bahwa $\int\!\int\abs f = +\infty$: jadi hipotesis
keterintegralan Fubini bukanlah hiasan.
\end{exercise}

\begin{exercise}[$\star\star$]\label{exo:b3:product:5}
(a) Hitunglah $\mathbf 1_{\intcc01} * \mathbf 1_{\intcc01}$
secara eksplisit (yakni sebuah fungsi tenda), lalu bentuk umum $(\mathbf 1 * \mathbf 1 *
\mathbf 1)$.
(b) Tunjukkan bahwa $\operatorname{supp}(f * g) \subseteq
\overline{\operatorname{supp}f + \operatorname{supp}g}$.
(c) Tunjukkan bahwa jika $f \in L^1$ dan $g$ terbatas serta \omterm{def:b3:topology:continuity}{kontinu}, maka
$f * g$ \omterm{def:b3:topology:continuity}{kontinu}. \emph{(Pakai kekonvergenan terdominasi lewat \omterm{def:b3:topology:continuity}{kekontinuan}
translasi pada $g$ yang terbatas itu.)}
\end{exercise}

\begin{exercise}[$\star\star$]\label{exo:b3:product:6}
(a) Tunjukkan bahwa simpleks $\Delta_d = \{x \in \intco0\infty^d
: x_1 + \dots + x_d \leq 1\}$ bervolume $\frac1{d!}$
\emph{(lewat induksi dan Fubini)}.
(b) Pulihkan $v_2 = \pi$ dan $v_3 = \frac{4\pi}3$ dari
\cref{thm:b3:product:ballvolume}, lalu tunjukkan bahwa
$\lambda_d(\text{elipsoid bersetengah sumbu } a_i) =
v_d\prod a_i$.
\end{exercise}

\begin{exercise}[$\star\star$]\label{exo:b3:product:7}
Untuk $s > 0$ yang mana integral berikut berhingga? Benarkan dengan \omterm{ex:b3:product:polar}{koordinat
kutub}:
\[
\int_{B(0,1)\subseteq\R^2}\frac{\dd x\,\dd y}{(x^2 +
y^2)^{s}},
\qquad
\int_{\R^2\setminus B(0,1)}\frac{\dd x\,\dd y}{(x^2 +
y^2)^{s}} .
\]
Perumumkanlah ke $\R^d$ (dengan ambang $s < d/2$ dan $s > d/2$).
\end{exercise}

\begin{exercise}[$\star\star\star$]\label{exo:b3:product:8}
(Beta--Gamma) Untuk $p, q > 0$, misalkan $B(p, q) = \int_0^1t^{p-1}(1
- t)^{q-1}\dd t$. Mulailah dari $\Gamma(p)\Gamma(q)$ sebagai
integral rangkap, lalu substitusikan $(x, y) = (uv,\, u(1 - v))$ (yakni sebuah
difeomorfisma dari kuadran terbuka ke
$\intoo0\infty\times\intoo01$; hitunglah Jacobinya $= u$) lalu
simpulkan\index{fungsi Beta}
\[
B(p, q) = \frac{\Gamma(p)\,\Gamma(q)}{\Gamma(p + q)} .
\]
Turunkan $\int_0^{\pi/2}\sin^{2p-1}\theta\cos^{2q-1}\theta\,
\dd\theta = \frac12B(p,q)$ beserta nilai integral
Wallis $W_n = \int_0^{\pi/2}\sin^n$.
\end{exercise}

\begin{exercise}[$\star\star$]\label{exo:b3:product:9}
(Rumus pemindahan) Misalkan $T \colon (X, \mathcal A, \mu) \to (Y,
\mathcal B)$ \omterm{def:b3:lebesgue:measurable}{terukur} dan $T_*\mu(B) = \mu(T^{-1}(B))$
yakni \emph{ukuran peta maju}\index{ukuran peta maju}.
Tunjukkan bahwa untuk setiap $g \geq 0$ yang \omterm{def:b3:lebesgue:measurable}{terukur} pada $Y$:
\[
\int_Y g\,\dd(T_*\mu) = \int_X g\circ T\,\dd\mu
\]
(lewat mesin bakunya). Lalu bandingkan dengan
\cref{thm:b3:product:linearchange}: keterangan tambahan apa yang dibawa
rumus penggantian variabelnya, yang tak dibawa
rumus pemindahan abstraknya? \emph{(Sebab rumus pemindahannya tak pernah
mengenali $T_*\mu$; sedangkan teorema penggantian variabelnya menghitung
$T_*\lambda_d$ secara eksplisit sebagai \omterm{def:b3:measure:measure}{ukuran} berkepadatan.)}
\end{exercise}

\begin{exercise}[$\star\star\star$]\label{exo:b3:product:10}
(Momen Gauss) Dengan \omterm{ex:b3:product:polar}{koordinat kutub} dan Fubini, hitunglah
untuk bobot Gauss baku pada $\R^d$:
\[
\int_{\R^d}\norm x_2^2\;\eu^{-\norm x_2^2}\,\dd x
\qquad\text{dan}\qquad
\int_{\R^d}x_1^2\,\eu^{-\norm x^2_2}\,\dd x,
\]
lalu periksa kekonsistenannya (sebab $\norm x^2 = \sum x_i^2$), dan turunkan
momen kedua \omterm{def:b3:measure:measure}{ukuran}
$\pi^{-d/2}\eu^{-\norm x^2}\dd x$.
\end{exercise}

\begin{exercise}[$\star\star$]\label{exo:b3:product:11}
(Grafik dan hipograf) Misalkan $f \colon \R^d \to \intco0\infty$
\omterm{def:b3:lebesgue:measurable}{terukur}.
(a) Tunjukkan bahwa \emph{hipograf} $H = \{(x, y) \in
\R^d\times\R : 0 < y < f(x)\}$ \omterm{def:b3:lebesgue:measurable}{terukur} di
$\R^{d+1}$ dengan
\[
\lambda_{d+1}(H) = \int_{\R^d}f\,\dd\lambda_d :
\]
yakni ``integralnya adalah luas di bawah grafiknya'', akhirnya menjadi sebuah
teorema. \emph{(Pakai penampangnya; lalu Tonelli.)}
(b) Tunjukkan bahwa grafik $\{(x, f(x)) : x \in \R^d\}$ merupakan
himpunan nol di $\R^{d+1}$.
(c) Turunkan bukti dua baris bahwa bola $S^{d-1}$ bersifat
nol-Lebesgue di $\R^d$.
\end{exercise}

\begin{exercise}[$\star\star$]\label{exo:b3:product:12}
(Sebuah integral rangkap yang termasyhur) Dengan deret geometri dan
Tonelli pada $\intoo01^2$, buktikan bahwa
\[
\int_0^1\!\!\int_0^1\frac{\dd x\,\dd y}{1 - xy}
= \sum_{n\geq1}\frac1{n^2} = \zeta(2),
\qquad
\int_0^1\!\!\int_0^1\frac{\dd x\,\dd y}{1 + xy}
= \sum_{n\geq1}\frac{(-1)^{n-1}}{n^2} = \frac{\zeta(2)}2 .
\]
(Untuk identitas deret keduanya: pisahkan indeks genap dan gasalnya.)
Dengan $\zeta(2) = \frac{\pi^2}6$ (\cref{ch:b3:spectral}), dua
integral yang tampak tak berbahaya bernilai $\frac{\pi^2}6$ dan
$\frac{\pi^2}{12}$; lalu di manakah tepatnya hipotesis kepositifan Tonelli
menjalankan tugasnya?
\end{exercise}

\section{Soal: rumus Stirling}

\begin{problem}[{Soal akhir pekan --- $n! \sim
\sqrt{2\pi n}\,(n/\eu)^n$, lewat kekonvergenan terdominasi}]
\label{pb:b3:product:1}
Rumus Stirling menguasai setiap pencacahan asimtotik di buku ini
--- yakni volume bola, koefisien binomial, dan bentuk lokal teorema
limit pusat. Kita membuktikannya dari integral $\Gamma$
(\cref{ex:b3:lebesgue:gamma}) dengan \emph{metode Laplace},
dalam bentuk kekonvergenan terdominasinya yang paling bersih, lalu memungut
buah hasilnya.\index{rumus Stirling}

\medskip
\noindent\textbf{Bagian I --- Rumusnya.} Untuk $t > 0$,
$\Gamma(t + 1) = \int_0^\infty x^{t}\eu^{-x}\dd x$.
\begin{enumerate}
  \item Substitusikan $x = t + \sqrt t\,u$ lalu tunjukkan bahwa
        \[
        \frac{\Gamma(t+1)}{t^{t}\eu^{-t}\sqrt t}
        = \int_{-\sqrt t}^{+\infty}
        \exp\Bigl(t\ln\Bigl(1 + \frac u{\sqrt
        t}\Bigr) - \sqrt t\,u\Bigr)\,\dd u
        \;=\;\int_\R g_t(u)\,\dd u,
        \]
        dengan $g_t(u) = \exp\bigl(t\ln(1 + u/\sqrt t) -
        \sqrt t\,u\bigr)\mathbf 1_{u > -\sqrt t}$.
  \item Tunjukkan limit titik demi titiknya: untuk setiap $u$ yang tetap, $g_t(u)
        \to \eu^{-u^2/2}$ saat $t \to +\infty$ \emph{(uraikan
        $\ln(1 + h)$ sampai orde kedua)}.
  \item Dominasinya. Misalkan $\varphi(h) = \ln(1 + h) - h$, sehingga
        $g_t(u) = \exp\bigl(t\,\varphi(u/\sqrt t)\bigr)$ untuk
        $u > -\sqrt t$. Buktikan kedua batas
        \[
        \varphi(h) \leq -\frac{h^2}4 \quad (-1 < h \leq 1),
        \qquad
        \varphi(h) \leq -c\,h \quad (h \geq 1),\ \ c = 1 -
        \ln 2 > 0
        \]
        \emph{(pelajarilah $\varphi(h) + \frac{h^2}4$ dan $\varphi(h)
        + ch$: hitunglah turunannya lalu periksa tandanya pada
        masing-masing selangnya)}. Lalu turunkan, untuk $t \geq 1$:
        \[
        g_t(u) \leq \eu^{-u^2/4}\ \ (\abs u \leq \sqrt t),
        \qquad
        g_t(u) \leq \eu^{-cu}\ \ (u \geq \sqrt t),
        \]
        sehingga $g_t(u) \leq \eu^{-u^2/4} +
        \eu^{-cu}\,\mathbf 1_{u > 0}$: yakni pendominasi \omterm{def:b3:lebesgue:l1}{terintegralkan}
        yang tak bergantung pada $t \geq 1$.
  \item Simpulkan dengan teorema kekonvergenan terdominasi dan
        integral Gauss (\cref{ex:b3:product:polar}):
        \[
        \Gamma(t + 1) \;\sim\;
        \sqrt{2\pi t}\;\Bigl(\frac t\eu\Bigr)^{t}
        \qquad (t \to +\infty),
        \]
        dan khususnya $n! \sim \sqrt{2\pi
        n}\,(n/\eu)^n$.
\end{enumerate}

\medskip
\noindent\textbf{Bagian II --- Buah hasilnya.}
\begin{enumerate}[resume]
  \item (Wallis) Dari \cref{exo:b3:product:8}, rumus bertipe $W_{2n} =
        \frac\pi2\binom{2n}n4^{-n}$: turunkanlah
        $\binom{2n}{n} \sim \frac{4^n}{\sqrt{\pi n}}$ dari
        Stirling, lalu periksalah terhadap rekursi $W_{n} =
        \frac{n-1}nW_{n-2}$.
  \item (Volume bolanya runtuh) Tunjukkan bahwa
        \[
        v_d = \frac{\pi^{d/2}}{\Gamma(\frac d2 + 1)}
        \;\sim\; \frac{1}{\sqrt{\pi d}}
        \Bigl(\frac{2\pi\eu}{d}\Bigr)^{d/2},
        \]
        sehingga $v_d \to 0$ lebih cepat daripada barisan geometri mana pun;
        lalu carilah dimensi yang memaksimumkan $v_d$ (secara numerik: $d =
        5$).
  \item (Pemusatan binomialnya --- yakni cuplikan
        \cref{ch:b3:clt}) Dengan Stirling, tunjukkan taksiran
        lokalnya, untuk $k = n/2 + s\sqrt n/2$ dengan $s$ yang tetap
        dan $n$ genap:
        \[
        2^{-n}\binom{n}{k} \;\sim\;
        \sqrt{\frac{2}{\pi n}}\;\eu^{-s^2/2},
        \]
        yakni profil Gauss diskretnya: de Moivre--Laplace dalam
        embrio.
  \item Di manakah tepatnya bukti Bagian I memakai: (i) kekonvergenan monoton atau
        terdominasi; (ii) integral Gaussnya; (iii) sifat keinvarianan
        \omterm{def:b3:measure:lebesgueouter}{ukuran Lebesgue}? Satu kalimat masing-masing.
\end{enumerate}

\medskip
\noindent\textbf{Bagian III --- Suku galatnya: Stirling dengan
palang.} Tetapkan $d_n = \ln n! - \bigl(n + \tfrac12\bigr)\ln n + n -
\ln\sqrt{2\pi}$, sehingga Bagian I mengatakan $d_n \to 0$.
\begin{enumerate}[resume]
  \item Tunjukkan bahwa $d_n - d_{n+1} = \bigl(n +
        \tfrac12\bigr)\ln\bigl(1 + \tfrac1n\bigr) - 1$.
  \item Dengan $t = \frac1{2n+1}$, periksalah $\frac{n+1}n =
        \frac{1+t}{1-t}$ lalu uraikan:
        \[
        d_n - d_{n+1} = \frac{t^2}3 + \frac{t^4}5 +
        \frac{t^6}7 + \cdots,
        \]
        lalu turunkan batas dua sisinya
        \[
        \frac1{3(2n+1)^2} \;<\; d_n - d_{n+1} \;<\;
        \frac1{12n} - \frac1{12(n+1)} .
        \]
  \item Teleskopkanlah (dengan memakai $d_m \to 0$) lalu periksa identitas
        aljabar yang menyenangkan $\frac1{3(2m+1)^2} >
        \frac1{12m+1} - \frac1{12(m+1)+1}$ untuk $m \geq 1$, untuk
        memperoleh pengapitan klasiknya
        \[
        \sqrt{2\pi n}\Bigl(\frac n\eu\Bigr)^n
        \eu^{1/(12n+1)} \;<\; n! \;<\;
        \sqrt{2\pi n}\Bigl(\frac n\eu\Bigr)^n\eu^{1/(12n)} .
        \]
  \item Dua akibatnya: (a) galat nisbi rumus Stirling
        bernilai $< 10^{-6}$ segera setelah $n \geq 83\,334$;
        (b) taksirlah $100!$ sampai empat angka berarti dengan tangan
        dari pengapitannya (dengan $100! \approx 9.3326\cdot
        10^{157}$), lalu takjublah sejenak pada ketelitian sebuah
        rumus asimtotik pada $n$ yang sangat berhingga.
\end{enumerate}

\medskip
\noindent\textbf{Bagian IV --- Jalur Wallis: Stirling
tanpa fungsi Gauss.} Secara historis konstanta $\sqrt{2\pi}$
datang dari Wallis, bukan dari Gauss; dan bagian ini membuktikan kembali Stirling
secara bebas dari Bagian I--II, sehingga membuktikan kembali
integral Gaussnya. Misalkan $W_n = \int_0^{\pi/2}\sin^n\theta\,
\dd\theta$.
\begin{enumerate}[resume]
  \item Tegakkanlah $W_n = \frac{n-1}nW_{n-2}$ (lewat pengintegralan
        parsial), bentuk tertutupnya
        \[
        W_{2n} = \frac\pi2\binom{2n}n4^{-n}, \qquad
        W_{2n+1} = \frac{4^n}{(2n+1)\binom{2n}n},
        \]
        beserta identitas $W_nW_{n-1} = \frac\pi{2n}$.
  \item Dari kemonotonan $(W_n)$ turunkan
        $W_{2n}/W_{2n+1} \to 1$, lalu
        \[
        W_{2n} \sim \frac12\sqrt{\frac\pi n}
        \qquad\text{dan}\qquad
        \binom{2n}n4^{-n}\sqrt n \longrightarrow
        \frac1{\sqrt\pi} :
        \]
        yakni teorema Wallis, yang diperoleh tanpa Stirling.
  \item Tunjukkan, lewat peneleskopan Bagian III saja (tanpa perlu nilai
        konstantanya), bahwa $e_n = \ln n! - (n +
        \frac12)\ln n + n$ konvergen ke suatu limit $\ell$;
        yang setara dengan $n! \sim K\,n^{n+1/2}\eu^{-n}$ dengan $K =
        \eu^\ell > 0$ yang belum dikenali.
  \item Sisipkan keasimtotikan ini ke $\binom{2n}n4^{-n}\sqrt n$
        lalu kenali, dengan memakai pertanyaan 14, satu-satunya nilai
        yang mungkin: $K = \sqrt{2\pi}$. Rakitlah logikanya: Bagian
        III--IV bersama memberikan bukti kedua yang \omterm{def:b3:complete:complete}{lengkap} bagi
        Stirling --- dan karena itu, dengan menjalankan substitusi Bagian I
        secara mundur, sebuah penilaian bebas atas
        $\int_\R\eu^{-u^2/2}\dd u = \sqrt{2\pi}$. Yakni dua tiang,
        yang masing-masingnya menopang yang lain.
\end{enumerate}

\medskip
\noindent\textbf{Bagian V --- Buah hasil terakhirnya.}
\begin{enumerate}[resume]
  \item (Profil lokal yang \omterm{def:b3:complete:complete}{lengkap}) Untuk bilangan bulat $\abs j \leq
        K\sqrt n$ (dengan $K$ yang tetap), tunjukkan bahwa
        \[
        \frac{\binom{2n}{n+j}}{\binom{2n}{n}} =
        \prod_{i=1}^{\abs j}\frac{n - i + 1}{n + i}
        = \exp\Bigl(-\frac{j^2}n +
        O\Bigl(\frac1{\sqrt n}\Bigr)\Bigr),
        \]
        secara seragam terhadap $j$ \emph{(ambillah logaritmanya lalu pakailah $\ln
        \frac{1-x}{1+y} = -(x + y) + O(x^2 + y^2)$)}. Inilah
        versi dua sisi pertanyaan 7 dan taksiran persis
        yang dikutip soal akhir pekan \cref{ch:b3:clt}.
  \item (Cuplikan Poisson) Tunjukkan dengan Stirling bahwa
        $\eu^{-n}\dfrac{n^n}{n!} \sim \dfrac1{\sqrt{2\pi n}}$:
        sehingga modus hukum Poisson bermean besar $n$ membawa
        massa $\approx (2\pi n)^{-1/2}$, persis seperti yang akan diramalkan
        teorema limit pusatnya.
  \item (Nisbah Gamma) Untuk $a \in \intoo01$, buktikan
        $\dfrac{\Gamma(n + a)}{\Gamma(n)\,n^a} \to 1$ dengan memakai
        batas kemiringan kelog-cembungan pada
        \cref{pb:b3:lebesgue:1} (yakni pertanyaan 14 di sana), lalu
        perluas ke setiap real $a > 0$ lewat persamaan
        fungsionalnya. (Inilah arti ``$\Gamma(t+1) \sim$
        Stirling'' di antara bilangan bulatnya.)
  \item (Bola, sekali lagi) Dari $v_d = \pi^{d/2}/\Gamma(\frac d2
        + 1)$: tabelkanlah $v_1, \dots, v_7$ secara persis, periksalah
        keunimodalannya lewat $\frac{v_d}{v_{d-2}} = \frac{2\pi}d$
        (yang naik selama $d < 2\pi$ dan turun sesudahnya), lalu
        buktikan identitas pembangkit yang mencolok
        \[
        \sum_{k\geq0}v_{2k}\,x^{2k} = \eu^{\pi x^2} :
        \]
        yakni semua \omterm{thm:b3:product:ballvolume}{volume bola satuan} berdimensi genap yang terkemas dalam satu
        eksponensial.
  \item (Keasimtotikan entropi) Untuk $\alpha \in \intoo01$ yang tetap
        dengan $\alpha n \in \N$, turunkan dari Stirling
        \[
        \binom{n}{\alpha n} \;\sim\;
        \frac{\eu^{n\,H(\alpha)}}
        {\sqrt{2\pi\,\alpha(1-\alpha)\,n}},
        \qquad
        H(\alpha) = -\alpha\ln\alpha -
        (1-\alpha)\ln(1-\alpha) :
        \]
        yakni laju pertumbuhan eksponensial koefisien binomialnya adalah
        \emph{entropi} $H$ --- lalu periksalah bahwa $\alpha =
        \frac12$ memulihkan pertanyaan 5, dan bahwa $H(\alpha) <
        \ln2$ untuk $\alpha \neq \frac12$ (sehingga binomial yang tak berpusat
        bersifat eksponensial terabaikan di dalam $2^n$).
  \item (Luas permukaan) Luas bola satuan $S^{d-1}$
        adalah $s_{d-1} = d\,v_d$ (yang dibuktikan sebagai
        \cref{exo:b3:forms:6} pada bab bentuk diferensial;
        di sini, terimalah sebagai definisinya). Tabelkanlah
        $s_0, \dots, s_6$, temukan yang terbesar (yakni $d - 1 =
        6$ dengan $s_6 = \frac{16\pi^3}{15} \approx 33.07$), lalu
        tunjukkan bahwa $s_{d-1} \to 0$ secara adigeometri pula ---
        sebab bola berdimensi tinggi, menurut setiap tolok ukur Euclid,
        bersifat lenyap kecil.
  \item (Suku koreksi pertamanya) Turunkan dari
        pengapitan pertanyaan 11 bahwa $d_n = \frac1{12n} +
        O\bigl(\frac1{n^2}\bigr)$, sehingga
        \[
        n! = \sqrt{2\pi n}\,\Bigl(\frac
        n\eu\Bigr)^{n}\Bigl(1 + \frac1{12n} +
        O\Bigl(\frac1{n^2}\Bigr)\Bigr).
        \]
        Periksalah di $n = 10$: rumus telanjangnya memberikan
        $3\,598\,696$ (dengan galat nisbi $8.3\cdot10^{-3}$),
        sedangkan yang terkoreksi $3\,628\,685$ terhadap $10! =
        3\,628\,800$ (dengan galat nisbi $3.2\cdot10^{-5}$) ---
        jadi satu suku deretnya membeli dua setengah angka.
  \item (Median $\Gamma$) Tunjukkan bahwa
        \[
        \frac{1}{\Gamma(t+1)}
        \int_0^{t} x^{t}\eu^{-x}\,\dd x
        \;\longrightarrow\; \frac12
        \qquad (t \to +\infty) :
        \]
        yakni secara asimtotik, tepat separuh massa
        integran $\Gamma$ terletak di bawah modusnya $x = t$.
        \emph{(Jalankan substitusi Bagian I pada integral yang
        terpancung; sebab pendominasi pertanyaan 3 sudah
        tersedia.)}
  \item (Entropi, secara tak asimtotik) Untuk $\alpha \in
        \intoc0{\frac12}$ buktikan batas yang berlaku bagi
        \emph{setiap} $n \geq 1$:
        \[
        \sum_{k=0}^{\lfloor\alpha n\rfloor}\binom nk
        \;\leq\; \eu^{n\,H(\alpha)} ,
        \]
        dengan membandingkan jumlahnya terhadap $\sum_k\binom
        nk\lambda^{k-\alpha n}$ untuk kemiringan $\lambda =
        \frac{\alpha}{1-\alpha} \leq 1$. Periksalah bahwa pemilihan
        $\lambda$ ini optimal, lalu damaikan dengan
        pertanyaan 21: sebab laju eksponensial $H(\alpha)$ pada
        pernyataan asimtotiknya tercapai oleh sebuah
        ketaksamaan satu baris tanpa keasimtotikan sama sekali.

\end{enumerate}
\end{problem}
